\documentclass[french]{book}

\usepackage[utf8]{inputenc}
\usepackage[T1]{fontenc}
\usepackage{lmodern}
\usepackage{geometry}
\usepackage{graphicx}
\usepackage{babel}
\usepackage{amsmath}
\usepackage{amsthm}
\usepackage{amssymb}
\usepackage{hyperref}
\usepackage{stmaryrd}
\usepackage{enumitem}
\usepackage{tcolorbox}
\usepackage{dsfont}
\usepackage{multirow}
\usepackage{etoc}
\usepackage{titlesec}
\usepackage{esint}
\usepackage{cancel}
\usepackage{mathdots}


\setlist[itemize]{label=$\bullet$}


\renewcommand{\P}{\mathbb{P}}
\newcommand{\N}{\mathbb{N}}
\newcommand{\Z}{\mathbb{Z}}
\newcommand{\Q}{\mathbb{Q}}
\newcommand{\R}{\mathbb{R}}
\newcommand{\C}{\mathbb{C}}
\newcommand{\K}{\mathbb{K}}
\renewcommand{\L}{\mathbb{L}}
\newcommand{\U}{\mathbb{U}}
\newcommand{\st}{^*}
\newcommand{\tq}{\middle|}
\newcommand{\inv}{^{-1}}
\newcommand{\E}{\mathbb{E}}
\newcommand{\F}{\mathbb{F}}
\newcommand{\V}{\mathbb{V}}
\renewcommand{\ker}{\text{Ker}}
\newcommand{\im}{\text{Im}}
\newcommand{\card}{\text{Card}}
\newcommand{\Tr}{\text{Tr}}

\theoremstyle{definition}
\newtheorem{ex}{Exercice}

\theoremstyle{theorem}
\newtheorem*{ind}{Indication}

\title{Exercices de première année}
\begin{document}
\maketitle
\chapter{Groupe symétrique}
\begin{ex}[Générateurs de $\mathcal{S}_n$]
	Soit $n\geqslant 2$.
	\begin{enumerate}
		\item Montrer que toute permutation de $\llbracket1,n\rrbracket$ peut s'écrire comme produit de transpositions de la forme $\begin{pmatrix}
			1&i
		\end{pmatrix}$ pour $2\leqslant i\leqslant n$.
		\item Montrer que toute permutation de $\llbracket1,n\rrbracket$ peut s'écrire comme produit de transpositions de la forme $\begin{pmatrix}
			i&i+1
		\end{pmatrix}$.
		\item Montrer que toute permutation de $\llbracket1,n\rrbracket$ peut s'écrire comme produit des cycles $\begin{pmatrix}
			1&2
		\end{pmatrix}$ et $\begin{pmatrix}
		1&2&\cdots&n
		\end{pmatrix}$.
	\end{enumerate}
\end{ex}
\begin{ex}[Signature d'une restriction]
	Soit $\sigma\in\mathcal{S}_n$ et $p\in\llbracket1,n\rrbracket$. On suppose que le support de $\sigma$ est inclus dans $\llbracket 1,p\rrbracket$ et on pose $\hat\sigma=\displaystyle \sigma_{\mid\llbracket1,p\rrbracket}^{\mid\llbracket1,p\rrbracket}$.
	\begin{enumerate}
		\item Montrer que $\hat{\sigma}$ est bien définie et qu'elle est un élément de $\mathcal{S}_p$.
		\item Montrer que $\hat{\sigma}$ et $\sigma$ ont la même signature.
	\end{enumerate}
\end{ex}
\begin{ex}
	$2022$ sont rangés de gauche à droite. Parmi celles qui sont proposées ci-dessous, quelles sont les transformations qui, itérées, permettent de ranger les objets dans l'ordre inverse ?
	\begin{enumerate}[label=\roman*)]
		\item On peut permuter entre eux deux objets arbitraires.
		\item On peut permuter entre eux deux objets arbitraires voisins.
		\item On peut permuter entre eux deux objets arbitraires, à condition qu'ils soient séparés par exactement un objet.
		\item On peut permuter entre eux deux objets arbitraires, à condition qu'ils soient séparés par exactement deux objets.
		\item On peut choisir arbitrairement trois objets et leur faire subir une permutation circulaire.
		\item On peut choisir quatre objets arbitraires et les ranger sur les places qu'ils occupaient précédemment, mais dans l'ordre inverse.
	\end{enumerate}
\end{ex}
\begin{ex}
	Calculer la signature de la permutation suivante : $$\sigma=\begin{pmatrix}
		1&2&\cdots&n&n+1&n+2&\cdots&2n\\
		1&3&\cdots&2n-1&2&4é\cdots&2n
	\end{pmatrix}$$
\end{ex}
\begin{ex}
	\begin{enumerate}
		\item Pour quelles valeurs de $n\in\N\st$ le groupe $(\mathcal{S}_n,\circ)$ est-il commutatif ?
		\item $\star$ Déterminer le centre $Z(\mathcal{S}_n)$ de $\mathcal{S}_n$ en fonction de $n$.
	\end{enumerate}
\end{ex}
\begin{ex}
	Dans $\mathcal{S}_n$, trouver toutes les permutations qui commutent avec le cycle $\begin{pmatrix}
		1&2&\cdots&n
	\end{pmatrix}$.
\end{ex}
\begin{ex}[Conjugaison]
	On dit que deux permutations $\sigma$ et $\sigma'$ de $\mathcal{S}_n$ sont conjuguées lorsque : $$\exists\tau\in\mathcal{S}_n,\ \sigma'=\tau\inv\sigma\tau$$
	\begin{enumerate}
		\item Montrer que la relation de conjugaison est une relation d'équivalence?
		\item Soit $(a\ b)$ une transposition et $\tau\in\mathcal{S}_n$. Montrer que $\tau(a\ b)\tau\in=(\tau(a)\ \tau(b))$. En déduire que deux transpositions sont toujours conjuguées.
		\item Généraliser à deux $p$-cycles.
	\end{enumerate}
\end{ex}
\begin{ex}
	Soit $n$ un entier supérieur ou égal à $2$. Montrer que la signature est l'unique morphisme surjectif de $\mathcal{S}_n$ dans $\{-1,1\}$.
\end{ex}
\begin{ex}
	Soit $n\geqslant 1$. Trouver tous les morphismes de groupe de $(\mathcal{S}_n,\circ)$ dans $(\C\st,\times)$.
\end{ex}
\begin{ex}
	Soit $\sigma\in\mathcal{S}_n$ que l'on décompose en produit de cycles à supports disjoints : $\sigma=c_1\cdots c_p$. On note $n_i$ la longueur du cycle $c_i$. Montrer que l'ordre de $\sigma$ dans $\mathcal{S}_n$ vaut $$o(\sigma)=n_1\vee\cdots\vee n_p$$
\end{ex}
\begin{ex}[$\star$ Le jeu du Taquin]
	Voici une configuration initiale $$\begin{pmatrix}
		1&2&3&4\\
		5&6&7&8\\
		9&10&11&12\\
		13&14&15&\emptyset
	\end{pmatrix}$$ On souhaite, en modifiant le tableau comme si on faisait glisser les numéros, passer à la configuration suivante : $$\begin{pmatrix}
	1&2&3&4\\
	5&6&7&8\\
	9&10&11&12\\
	13&15&14&\emptyset
	\end{pmatrix}$$ Qu'en pensez-vous ?
\end{ex}
\begin{ex}[$\star\star$ Automorphismes intérieurs de $\mathcal{S}_n$ et conjugaison]
	On rappelle que deux permutations $\sigma$ et $\sigma'$ de $\mathcal{S}_n$ sont conjuguées lorsque : $$\exists\tau\in\mathcal{S}_n,\ \sigma'=\tau\inv\sigma\tau$$
	\begin{enumerate}
		\item Soit $\tau\in\mathcal{S}_n$. Vérifier que l'application $\varphi:\sigma\mapsto\tau\inv\sigma\tau$ est un automorphisme de $\mathcal{S}_n$. On dit que c'est un automorphisme intérieur.
		\item Soit $\varphi$ un automorphisme de $\mathcal{S}_n$. Montrer que $\varphi$ est un automorphisme intérieur si, et seulement s'il transforme toute transposition en une transposition.
	\end{enumerate}
\end{ex}
\begin{ex}[$\star\star$ IMO]
	Soit $p$ un nombre premier impair. Trouver le nombre de parties $A$ de $\llbracket 1,2p\rrbracket$ telles que : 
	\begin{itemize}
		\item $A$ est de cardinal $p$ ;
		\item la somme des éléments de $A$ est divisible par $p$.
	\end{itemize}
\end{ex}
\begin{ex}[$\star\star$ Sous-groupes alternés de $\mathcal{A}_n$ pour $n\geqslant 5$]
	On dit qu'un sous-groupe $H$ d'un groupe $G$ est distingué dans $G$ lorsque $\forall x\in G,\ xH=Hx$. Montrer que si $n\geqslant 5$, alors le groupe alterné $\mathcal{A}_n$ n'admet pas d'autres sous-groupes distingués de $\{\text{id}\}$ et $\mathcal{A}_n$.
\end{ex}


\chapter{Déterminants}
\begin{ex}
	Soit $A=(a_{i,j})_{1\leqslant i,j\leqslant n}\in\mathcal{M}_n(\K)$. On définit $A'=((-1)^{i+j}a_{i,j})_{1\leqslant i,j\leqslant n}\in\mathcal{M}_n(\K)$. Montrer que $\det(A')=\det(A)$.
\end{ex}
\begin{ex}
	Quel est le déterminant d'un projecteur ? Quel est le déterminant d'une symétrie ?
\end{ex}
\begin{ex}
	Soit $(a,b,c,d)\in\R^4$. Calculer sous forme factorisée les déterminants suivants.
	\begin{enumerate}
		\item $\displaystyle\begin{vmatrix}
			0&a&b\\
			a&0&c\\
			b&c&0
		\end{vmatrix}$
		\item $\displaystyle\begin{vmatrix}
			a&b&c\\
			c&a&b\\
			b&c&a
		\end{vmatrix}$
		\item $\displaystyle\begin{vmatrix}
			a+b&b+c&c+a\\
			a^2+b^2&b^2+c^2&c^2+a^2\\
			a^3+b^3&b^3+c^3&c^3+a^3
		\end{vmatrix}$
		\item $\displaystyle\begin{vmatrix}
			a&a&a&a\\
			a&b&b&b\\
			a&b&c&c\\
			a&b&c&d
		\end{vmatrix}$
		\item $\displaystyle\begin{vmatrix}
			a&c&c&b\\
			c&a&b&c\\
			c&b&a&c\\
			b&c&c&a
		\end{vmatrix}$
		\item $\displaystyle\begin{vmatrix}
		1&1&1\\
		\cos(a)&\cos(b)&\cos(c)\\
		\sin(a)&\sin(b)&\sin(c)
		\end{vmatrix}$
		\item $\displaystyle\begin{vmatrix}
			1&a&a^2&a^4\\
			1&b&b^2&b^4\\
			1&c&c^2&c^4\\
			1&d&d^2&d^4
		\end{vmatrix}$
	\end{enumerate}
\end{ex}
\begin{ex}
	Soit $M=(m_{i,j})_{1\leqslant i,j\leqslant n}\in\mathcal{M}_n(\R)$ telle que $$\begin{cases}
		\forall (i,j)\in\llbracket1,n\rrbracket^2,\ 0\leqslant m_{i,j}<1 \\
		\\
		\forall i\in\llbracket1,n\rrbracket,\ \displaystyle\sum_{j=1}^{n}m_{i,j}\leqslant 1
	\end{cases}$$ Montrer que $\lvert\det(M)\rvert<1$.
\end{ex}
\begin{ex}
	Soit $A=\displaystyle\begin{pmatrix}
		0&1&2\\
		1&1&1\\
		1&0&-1
	\end{pmatrix}$.
	\begin{enumerate}
		\item Soit $\lambda\in\C$. on dit que $\lambda$ est une valeur propre de la matrice $A$ lorsqu'il existe $X\in\C^3\setminus\{0\}$ tel que $AX=\lambda X$. Montrer que $\lambda$ est une valeur propre de $A$ si, et seulement si, $\det(\lambda I_3-A)=0$.
		\item Trouver une matrice $P\in\mathcal{GL}_3(\C)$ telle que $P\inv AP$ soit une matrice diagonale à déterminer.
		\item En déduire la valeur de $A^n$ pour $n\in\N\st$.
	\end{enumerate}
\end{ex}
\begin{ex}
	Soit $A\in\mathcal{M}_n(\K)$. Calculer $\text{rg}(\text{Com}(A))$ en fonction de $\text{rg}(A)$.
\end{ex}
\begin{ex}
	Trouver toutes les matrices $A\in\mathcal{M}_n(\K)$ telles que $$\forall B\in\mathcal{M}_n(\K),\ \det(A+B)=\det(A)+\det(B)$$
\end{ex}
\begin{ex}
	Soit $A\in\mathcal{M}_n(\K)$ une matrice de rang $1$. Montrer que $\det(I_n+A)=1+\Tr(A)$.
\end{ex}
\begin{ex}
	Soit $n\in\N\st$. Soit $J\in\mathcal{M}_n(\K)$ dont tous les coefficients valent $1$. Soit $A\in\mathcal{GL}_n(\K)$. On pose $B=A-J$ et on note $A\inv=(a_{i,j}')_{1\leqslant i,j\leqslant n}$. Montrer que $$\det(B)=\det(A)\left(1-\sum_{1\leqslant i,j\leqslant n}a_{i,j}'\right)$$
\end{ex}
\begin{ex}
	Soit $A\in\mathcal{M}_n(\R)$. Soit $H\in\mathcal{M}_n(\R)$ de rang $1$. Montrer l'inégalité $\det(A+H)\det(A-H)\leqslant\det(A)^2$.
\end{ex}
\begin{ex}
	Calculer $\displaystyle\begin{vmatrix}
		2&-1&&&&\\
		3&2&-1&&(0)&\\
		&\ddots&\ddots&\ddots&&\\
		&&\ddots&\ddots&\ddots&\\
		&(0)&&3&2&-1\\
		&&&&3&2
	\end{vmatrix}$.
\end{ex}
\begin{ex}
	Soit $(\lambda_1,\ldots,\lambda_n)\in\C^n$ une famille de scalaires deux à deux distincts. A l'aide d'un déterminant de Vandermonde, montrer que les fonctions $f:x\mapsto e^{\lambda_ix}$ forment une famille libre de $\C^\R$.
\end{ex}
\begin{ex}
	A l'aide d'un déterminant de Vandermonde, retrouver le résultat (de cours) selon lequel si deux polynômes $P$ et $Q$ deux $\K_n[X]$ coïncident en $n+1$ points distincts, alors ils sont égaux.
\end{ex}
\begin{ex}
	Soit $(a_0,\ldots,a_n)\in\R^{n+1}$ une famille de réels tous non nuls et deux à deux distincts. Montrer que la famille des applications : $$\varphi_i:P\mapsto\int_{0}^{a_i}P(t)\ dt$$ est une base de $\mathcal{L}(\R_n[X],\R)$.
\end{ex}
\begin{ex}
	Soit $(a_1,\ldots,a_n)\in\K^{n}$. On cherche à calculer le déterminant de Vandermonde $$V(a_0,\ldots,a_n)=\begin{vmatrix}
		1&a_1&\cdots&a_1^{n-1}\\
		\vdots&\vdots&&\vdots\\
		\vdots&\vdots&&\vdots\\
		1&a_n&\cdots&a_n^{n-1}
	\end{vmatrix}$$ On propose ici une démonstration par récurrence sur $n\in\N\st$.
	\begin{enumerate}
		\item Dans un premier temps, on suppose $n\geqslant 2$ et que les $a_i$ sont deux à deux distincts. On suppose aussi que l'hypothèse de récurrence est vraie au rang $n-1$. Introduisons la fonction $D:\K\to\K$ définie par $$\forall x\in\K,\ D(x)=\begin{vmatrix}
			1&a_1&\cdots&a_1^{n-1}\\
			\vdots&\vdots&&\vdots\\
			1&a_n&\cdots&a_{n-1}^{n-1}\\
			1&x&\cdots&x^{n-1}
		\end{vmatrix}$$ Montrer qu'il existe $P\in\K_{n-1}[X]$ tel que $D$ soit la fonction polynomiale associée à $P$, puis factoriser $P$.
		\item Conclure.
	\end{enumerate}
\end{ex}
\begin{ex}
	On se donne $n$ complexes $a_0,\ldots,a_n$ et l'on cherche à calculer le déterminant de la matrice $$A=\begin{pmatrix}
		a_0&a_1&\cdots&\cdots&a_{n-1}\\
		a_{n-1}&a_0&a_1&\cdots&a_{n-2}\\
		&\ddots&\ddots&\ddots&\\
		&&\ddots&\ddots&\ddots\\
		a_1&\cdots&\cdots&a_{n-1}&a_0
	\end{pmatrix} $$ On note $\omega=\exp\displaystyle\left(i\frac{2\pi}{n}\right)$, et on introduit la matrice $\Omega$ de terme général $\Omega_{i,j}\omega^{(i-1)(j-1)}$. On pose aussi $P=\displaystyle\sum_{k=0}^{n-1}a_kX^k\in\C[X]$.
	\begin{enumerate}
		\item Calculer $A\Omega$ en fonction de $P$ et $\omega$.
		\item Conclure.
	\end{enumerate} 
\end{ex}
\begin{ex}
	Soit $(a_1,\ldots,a_n)\in\K^n$. Calculer le déterminant de la matrice de terme générale $a_{\max(i,j)}$.
\end{ex}
\begin{ex}
	Calculer $$D_{n+1}=\begin{vmatrix}
		\binom{0}{0}&\binom{1}{1}&\cdots&\binom{n}{n}\\
		\binom{1}{0}&\binom{2}{1}&\cdots&\binom{n+1}{n}\\
		\vdots&\vdots&&\vdots\\
		\binom{n}{0}&\binom{n+1}{2}&\cdots&\binom{2n}{n}
	\end{vmatrix}$$
\end{ex}
\begin{ex}
	Soit $x\in\R$. Calculer le déterminant suivant, de taille $n\times n$ : $$D_n=\begin{vmatrix}
		2x&10&\cdots&0&0\\
		1&2x&1\cdots&0&0\\
		0&1&2x&\ddots&0&0\\
		\vdots&\vdots&\ddots&\ddots&\ddots&\vdots\\
		0&0&\cdots&1&2x&1\\
		0&0&\cdots&0&1&2x
	\end{vmatrix}$$ On pourra commencer par le cas où $x\in]-1,1[$ en posant $\theta=\arccos(x)$.
\end{ex}
\begin{ex}
	Soit $A$, $B$, $C$ et $D$ quatre matrices de $\mathcal{M}_n(\C)$ telles que $CD=CD$. On suppose de plus $D$ inversible. 
	\begin{enumerate}
	\item Montrer que $$\det\begin{pmatrix}
		A&B\\
		C&D
	\end{pmatrix}=\det(AD-BC)$$
	\item $\star$ Montrer que le résultat reste valable lorsque $D$ n'est plus nécessairement supposée inversible.
	\end{enumerate}
\end{ex}
\begin{ex}
	Soit $A$ et $B$ deux matrices de $\mathcal{M}_n(\R)$. Montrer que $$\det\begin{pmatrix}
		A&-B\\
		B&A
	\end{pmatrix}\geqslant0$$
\end{ex}
\begin{ex}
	On se donne deux matrices $A$ et $B$ de $\mathcal{M}_n(\R)$. On suppose que $A$ et $B$ sont semblables dans $\mathcal{M}_n(\C)$. Le but est de démontrer que $A$ et $B$ sont alors semblables dans $\mathcal{M}_n(\R)$.
	\begin{enumerate}
		\item Fixons $P\in\mathcal{GL}_n(\C)$ telle que $B=P\inv AP$. On pose $P_1=(\text{Re}(p_{i,j}))_{1\leqslant i,j\leqslant n}$ et $P_2=(\text{Im}(p_{i,j}))_{1\leqslant i,j\leqslant n}$. Montrer que $$\forall x\in\R,\ A(P_1+xP_2)=(P_1+xP_2)A$$
		\item Conclure.
	\end{enumerate}
\end{ex}
\begin{ex}
	Soit $(u_1,\ldots,u_p)$ une famille libre due $\Q$-espace vectoriel $\Q^n$. Montrer que, vue comme une famille du $\R$-espace vectoriel $\R^,$, cette famille est encore libre.
\end{ex}
\begin{ex}
	Soit $n\in\N\st$. Soit $P=X^n-\displaystyle\sum_{k=0}^{n-1}a_kX^k\in\mathcal{C}_n[X]$. On appelle matrice compagnon de $P$ la matrice $$M=\begin{pmatrix}
		0&&(0)&a_0\\
		1&\ddots&&a_1\\
		&\ddots&0&\vdots\\
		(0)&&1&a_{n-1}
	\end{pmatrix}\in\mathcal{M}_n(\C)$$
	\begin{enumerate}
		\item Montrer : $\forall \lambda\in\C,\ \det(\lambda I_n-M)=P(\lambda)$.
		\item Soit $u$ l'endomorphisme de $\C^n$ canoniquement associé à $M$ et $(e_i)_{1\leqslant i\leqslant n}$ la base canonique de $\C^n$. Exprimer $e_2,\ldots,e_n$ en fonction de $e_1$ et de $u$.
		\item Montrer que $P(M)=0$.
	\end{enumerate}
\end{ex}
\begin{ex}
	On note $\mathcal{GL}_n(\Z)=\left\{M\in\mathcal{GL}_n(\R)\cap\mathcal{M}_n(\Z)\ : \ M\inv\in\mathcal{M}_n(\Z)\right\}$. Soit $A\in\mathcal{M}_n(Z)$. Montrer l'équivalence : $$A\in\mathcal{GL}_n(\Z)\iff\det(A)=\pm 1$$
\end{ex}
\begin{ex}
	On définit $SL_n(\K)=\left\{M\in\mathcal{M}_n(\K)\ : \ \det(M)=1\right\}$. Montrer que $SL_n(\K)$ est un groupe engendré par les transvections sur les lignes.
\end{ex}
\begin{ex}
	\begin{enumerate}
		\item Soit $E$ un ensemble fini et $\varphi:E\to E$ une involution sans point fixe. Montrer qu'il existe une partie $F\subset E$ telle que $F\sqcup\varphi(F)=E$.
		\item Soit $A\in\mathcal{M}_{2n+1}(\Z)$ une matrice symétrique à coefficients entiers dont les termes diagonaux sont pairs. Montrer que $\det(A)$ est un entier pair.
	\end{enumerate}
\end{ex}
\begin{ex}[Résultant et discriminant]
	Dans tout l'exercice, on fixe $p$ et $q$ des entiers naturels tous deux non nuls. Soit $P=\displaystyle\sum_{k=0}^{p}a_kX^k$ et  $P=\displaystyle\sum_{l=0}^{q}b_lX^l$ dans $\C[X]$ avec $a_p\ne0$ et $b_q\ne0$. On définit le résultant de $P$ et $Q$ par $$\text{Res}(P,Q)=\begin{vmatrix}
	a_0&0&\cdots&0&b_0&0&\cdots&\cdots&\cdots&0 \\
	a_1&\ddots&\ddots&\vdots&b_1&\ddots&\ddots&&&\vdots\\
	\vdots&\ddots&\ddots&0&\vdots&\ddots&\ddots&\ddots&&\vdots\\
	\vdots&&\ddots&a_0&\vdots&&\ddots&\ddots&\ddots&\vdots\\
	\vdots&&&a_1&b_q&&&\ddots&\ddots&0\\
	\vdots&&&\vdots&0&\ddots&&&\ddots&b_0\\
	a_p&&&\vdots&\vdots&\ddots&\ddots&&&b_1\\
	0&\ddots&&\vdots&\vdots&&\ddots&\ddots&&\vdots\\
	\vdots&\ddots&\ddots&\vdots&\vdots&&&\ddots&\ddots&\vdots\\
	0&\cdots&0&a_p&0&\cdots&\cdots&\cdots&0&b_q
	\end{vmatrix}$$
	\begin{enumerate}
		\item En considérant l'application $$\begin{array}{lrcl}
			\varphi:&\C_{q-1}[X]\times\C_{p-1}[X]&\longrightarrow&\C_{p+q-1}[X]\\
			&(U,V)&\longmapsto &UP+VQ
		\end{array}$$ démontrer que $\text{Res}(P,Q)\ne 0$ si, et seulement si, $P\wedge Q=1$.
		\item On définit le discriminant de $P$ par $\Delta(P)=\displaystyle\frac{(-1)^{\frac{p(p-1)}{2}}}{a_p}\text{Res}(P,P')$. Qu'obtient-on pour un polynôme du second degré ?
		\item Montrer que $P$ admet une racine multiple si, et seulement si, sont discriminant est non nul.
		\item Application : déterminer une condition simple sur $(a,b)\in\C^2$ pour que le polynôme $X^"+aX+b$ possède une racine multiple.
		\item $\star$ On note $x_1,\ldots,z_p$ les racines complexes de $P$ comptées avec multiplicité. Exprimer $\Delta(P)$ sous forme factorisée en fonction des $z_i$. Retrouver le résultat de la question $3$.
	\end{enumerate}
\end{ex}
\begin{ex}[$\star$ Déterminant de Cauchy]
	Soit $(a_1,\ldots,a_n,b_1,\ldots,b_n)\in\K^{2n}$ tel que $\forall (i,j)\in\llbracket1,n\rrbracket^2,\ a_i+b_j\ne 0$.
	\begin{enumerate}
		\item Dans cette question, on suppose que les $b_j $ sont deux à deux distincts. Décomposer en éléments simples la fraction rationnelle $$F=\frac{(X-a_1)\cdots(X-a_{n-1})}{(X+b_1)\cdots(X+b_n)}$$
		\item Calculer le déterminant de la matrice $\displaystyle\left(\frac{1}{a_i+b_j}\right)_{1\leqslant i,j\leqslant n}$.
	\end{enumerate}
\end{ex}
\begin{ex}[$\star$ Produit tensoriel et formes $p$-linéaires]
	Soit $E$ un $\K$-espace vectoriel de dimension $n$, $(e_i)_{1\leqslant i\leqslant n}$ une base de $E$ et $p\in\N\st$.
	\begin{enumerate}
		\item Montrer que pour toute famille de scalaires $(\lambda_{i_1,\ldots,i_p})_{1\leqslant i_1,\ldots,i_p\leqslant n}$, il existe une unique forme $p$-linéaire $f$ telle que $$\forall (i_1,\ldots,i_p)\in\llbracket1,n\rrbracket^p,\ f(e_{i_1},\ldots,e_{i_p})=\lambda_{i_1,\ldots,i_p}$$
		\item En déduire que la famille $(e_{i_1}\st\otimes\cdots\otimes e_{i_p}\st)_{1\leqslant i_1,\ldots,i_p\leqslant n}$ est une base de l'espace des formes $p$-linéaires.
	\end{enumerate}
\end{ex}
\begin{ex}[$\star$ Comatrice du produit]
	Soit $A$ et $B$ deux matrices de $\mathcal{M}_n(\C)$. Montrer que la comatrice de $AB$ est égal au produit de la comatrice de $A$ par la comatrice de $B$.
\end{ex}
\begin{ex}[$\star$ Polynôme minimal et sur-corps]
	\begin{enumerate}
		\item Soit $A\in\mathcal{M}_n(\K)$ une matrice non nulle. Montrer qu'il existe un unique polynôme $\Pi_\K$ unitaire de degré minimal tel que $\Pi_\K(A)=0$.
		\item Soit $\L$ un sur-corps de $\K$, et $\Pi_\L$ le polynôme minimal de $A$, considérée comme une matrice de $\mathcal{M}_n(\L)$. Montrer que $\Pi_\K=\Pi_\L$.
	\end{enumerate}
\end{ex}
\begin{ex}[$\star\star$]
	Soit $n\in\N\st$ et $(f_1,\ldots,f_n)$ une famille libre de fonctions de $\R$ dans $\R$. Montrer qu'il existe $(x_1,\ldots,x_n)\in\R^n$ tel que $$\begin{vmatrix}
		f_1(x_1)&f_1(x_2)&\cdots&f_1(x_n)\\
		f_2(x_1)&f_2(x_2)&\cdots&f_2(x_n)\\
		\vdots&\vdots&&\vdots\\
		f_n(x_1)&f_n(x_2)&\cdots&f_n(x_n)
	\end{vmatrix}\ne 0$$
\end{ex}
\begin{ex}[$\star\star$ Les $2n+1$ cailloux]
	Soit $n\in\N\st$. On se donne $2n+1$ cailloux et on suppose que dès que l'on choisit un caillou au hasard, on peut toujours former deux tas de $n$ cailloux ayant le même poids. Montrer que tous les cailloux ont le même poids.
\end{ex}




\end{document}